\section{Analysis of Model-\texorpdfstring{$\CH$}{C2-H}}
\label{sec:model_c21}



Model-$\CH$ is the abandonment counterpart of Model-$\BH$, and the head-of-line
version of the class-$1$ abandonment model Model-$C_2$: the parameter setting is $\theta_1>0$ head-of-line with $\gamma_1=\gamma_2=\theta_2=0$. Only the \emph{customer at the head of Queue~1} abandons, at the fixed rate $\theta_1$ whenever $n_1\geq1$. As in Model-$\BH$, the constant rate $\theta_1\,\mathbf{1}_{\{n_1\ge1\}}$ contributes an \emph{algebraic} term to the fundamental equation rather than a derivative term, and the Kernel Method applies directly. Because an abandoning customer leaves the system altogether, the conservation of the total number of customers ---which held in every model without abandonment--- no longer applies.

\textbf{Stability.} Because only the head-of-line customer from class-$1$ abandons, the aggregate abandonment rate is bounded by $\theta_1$ regardless of the queue length ---unlike in Model-$C_2$, where it grows proportionally to $n_1$ and stabilises the priority queue at any load. A non-trivial stability condition is therefore required. As shown in Corollary~\ref{cor:C21:pi0} below, the chain is positive recurrent if and only if $\pi_0>0$, that is
\begin{equation}
    (\mu+\theta_1)(\mu-\lambda) + \lambda_1\theta_1 \;>\; 0,
    \qquad \lambda=\lambda_1+\lambda_2.
    \label{eq:C21:stability}
\end{equation}
This is \emph{weaker} than $\rho<1$: if $\rho<1$ then $\mu-\lambda>0$ and the condition holds trivially, but it can also hold with $\rho\ge1$ provided $\theta_1$ is large enough. Moreover, Equation~\eqref{eq:C21:stability} implies $\mu+\theta_1>\lambda_1$, i.e.\ $\rho_C:=\lambda_1/(\mu+\theta_1)<1$.

\subsection{Balance equations}

The abandonment transition $(n_1,n_2)\to(n_1-1,n_2)$ occurs at the fixed rate $\theta_1$ whenever $n_1\ge1$. Since both a service completion and a head-of-line abandonment from $(n_1+1,n_2)$ lead to $(n_1,n_2)$, the two mechanisms enter the balance equations through the common in-flow coefficient $\mu+\theta_1$. On $S$,
\begin{align}
    \left[\lambda_1+\lambda_2+\mu+\theta_1\right]\pi(n_1,n_2)
    &= (\mu+\theta_1)\pi(n_1+1,n_2)
      +\lambda_1\pi(n_1-1,n_2) \notag\\
    &\qquad
      +\lambda_2\pi(n_1,n_2-1),
    &&n_1\ge1,\;n_2\ge1, \label{eq:C21:interior} \\
    \left[\lambda_1+\lambda_2+\mu+\theta_1\right]\pi(n_1,0)
    &= (\mu+\theta_1)\pi(n_1+1,0)
      +\lambda_1\pi(n_1-1,0),
    &&n_1\ge1,\;n_2=0, \nonumber \\
    \left[\lambda_1+\lambda_2+\mu\right]\pi(0,n_2)
    &= (\mu+\theta_1)\pi(1,n_2)
      +\lambda_2\pi(0,n_2-1) \notag\\
    &\qquad
      +\mu\,\pi(0,n_2+1),
    &&n_1=0,\;n_2\ge1, \nonumber \\
    \left[\lambda_1+\lambda_2+\mu\right]\pi(0,0)
    &= (\lambda_1+\lambda_2)\pi_0
      + (\mu+\theta_1)\pi(1,0)
      + \mu\,\pi(0,1),
    &&\nonumber \\
    (\lambda_1+\lambda_2)\pi_0
    &= \mu\,\pi(0,0).
    &&\label{eq:C21:idle}
\end{align}

\begin{remark}[The substitution $\mu\mapsto\mu+\theta_1$]\label{rem:C21:mushift}
Equations~\eqref{eq:C21:interior}--\eqref{eq:C21:idle} are precisely the Model-$A$ balance equations with the service rate $\mu$ replaced by $\mu+\theta_1$ in every transition out of a state with $n_1\ge1$; transitions out of states with $n_1=0$ keep the rate $\mu$, since no waiting class-$1$ customer is present to abandon. The same substitution acts on the coefficient of $P(x,y)$ in the fundamental equation below, but not on its boundary terms, and it does not carry over to the stability condition~\eqref{eq:C21:stability}, which instead follows from the busy-period identity of Theorem~\ref{thm:gen:boundary} (see the proof of Corollary~\ref{cor:C21:pi0}).
\end{remark}

\subsection{Fundamental PGF equation}
We introduce the following result:
\begin{lemma}[Fundamental equation for Model-$\CH$]\label{lem:C21:fundamental}
    Assume positive recurrence. Then the joint PGF $P(x,y)$ and the boundary function $P_y(y)=P(0,y)$ satisfy
    \begin{equation}
        \begin{split}
            &\left[(\lambda_1+\lambda_2+\mu)xy - \mu y - \lambda_1 x^2y - \lambda_2 xy^2
                   + \theta_1 y(x-1)\right]P(x,y) \\
            &\quad=\quad \left[\mu(x-y) + \theta_1 y(x-1)\right]P_y(y)
                   - \mu x(1-y)\,\pi(0,0).
        \end{split}
        \label{eq:C21:fundamental}
    \end{equation}
\end{lemma}

\begin{proof}
    As in Lemma~\ref{lem:B21:fundamental}, deleting every term that contains $\theta_1$ from Equations~\eqref{eq:C21:interior}--\eqref{eq:C21:idle} leaves exactly the Model-$A$ balance equations, which reproduce the Model-$A$ fundamental equation~\eqref{eq:A:fundamental} after multiplication by $x^{n_1}y^{n_2}$, summation, and multiplication by $xy$. It remains to collect the terms containing $\theta_1$. The $\theta_1$-term on the left-hand side appears in every state with $n_1\geq1$,
    \begin{equation*}
        \theta_1\sum_{n_1\ge1}\sum_{n_2\ge0}\pi(n_1,n_2)\,x^{n_1}y^{n_2}
        = \theta_1\bigl[P(x,y)-P_y(y)\bigr],
    \end{equation*}
    while the $\theta_1$-term on the right-hand side enters every state $(n_1,n_2)$ from $(n_1+1,n_2)$,
    \begin{equation*}
        \theta_1\sum_{n_1\ge0}\sum_{n_2\ge0}\pi(n_1+1,n_2)\,x^{n_1}y^{n_2}
        = \theta_1\,x^{-1}\sum_{m_1\ge1}\sum_{n_2\ge0}\pi(m_1,n_2)\,x^{m_1}y^{n_2}
        = \theta_1\,x^{-1}\bigl[P(x,y)-P_y(y)\bigr],
    \end{equation*}
    with $m_1=n_1+1$. The net contribution is
    $\theta_1\bigl(1-x^{-1}\bigr)\bigl[P(x,y)-P_y(y)\bigr]$; multiplying by $xy$ gives $\theta_1 y(x-1)\bigl[P(x,y)-P_y(y)\bigr]$. Adding this to the Model-$A$ equation and rearranging yields Equation~\eqref{eq:C21:fundamental}.
\end{proof}

\subsection{Closed-form solution for the PGF}

The following theorem is the main result of this section; as for Model-$\BH$, its proof applies the Kernel Method of Model-$A$.

\begin{theorem}\label{thm:C21:PGF}
    Consider the two-class non-preemptive priority queue with Poisson arrival rates
    $\lambda_1,\lambda_2$, common exponential service rate $\mu$, and head-of-line class-$1$
    abandonment at rate $\theta_1>0$ (with $\gamma_1=\gamma_2=\theta_2=0$). Under the stability
    condition~\eqref{eq:C21:stability}, the joint PGF $P(x,y)$ is the
    rational function
    \begin{equation}
        P(x,y) \;=\;
        \frac{\mu(\mu+\theta_1)\,(1-y)\,\pi(0,0)}
        {\lambda_1\,\bigl(x^+(y)-x\bigr)
         \,\bigl[(\mu+\theta_1 y)\,x^*(y)-y(\mu+\theta_1)\bigr]},
        \label{eq:C21:PGF}
    \end{equation}
    where $x^*(y)$ and $x^+(y)$ are the two roots of the quadratic equation
    \begin{equation}
        \lambda_1 x^2 - \bigl[\lambda_1+\lambda_2(1-y)+\mu+\theta_1\bigr]x + (\mu+\theta_1) = 0,
        \label{eq:C21:quad}
    \end{equation}
    with $x^*(y)$ the root inside the unit disc, and $\pi(0,0)$ is given by
    Corollary~\ref{cor:C21:pi0}.
\end{theorem}
The theorem invokes the following Corollary:
\begin{corollary}\label{cor:C21:pi0}
    Under Equation~\eqref{eq:C21:stability}, the empty-state probabilities of Model-$\CH$ are
    \begin{equation}
        \pi_0 = \frac{(\mu+\theta_1)(\mu-\lambda)+\lambda_1\theta_1}
                     {\mu(\mu+\theta_1)+\lambda_1\theta_1},
        \qquad
        \pi(0,0) = \frac{\lambda\bigl[(\mu+\theta_1)(\mu-\lambda)+\lambda_1\theta_1\bigr]}
                        {\mu\bigl[\mu(\mu+\theta_1)+\lambda_1\theta_1\bigr]},
        \label{eq:C21:pi0_pi00}
    \end{equation}
    with $\lambda=\lambda_1+\lambda_2$. At $\theta_1\to0^+$ these reduce to $\pi_0\to1-\rho$ and
    $\pi(0,0)\to\rho(1-\rho)$, recovering Model-$A$.
\end{corollary}

As $\theta_1\to0^+$, Equation~\eqref{eq:C21:PGF} likewise recovers the Model-$A$ PGF of Theorem~\ref{thm:A:PGF} exactly. Theorem~\ref{thm:C21:PGF} also parallels Theorem~\ref{thm:B21:PGF} structurally: both PGFs are rational in $x$ with the pole governed by the outer kernel root $x^+(y)$, and the two quadratic kernels differ only in their constant term ---the $y$-dependent $\mu+\gamma_1y$ for jockeying versus the constant $\mu+\theta_1$ for abandonment. The remaining difference is the numerator input: for Model-$\BH$ conservation fixes $\pi(0,0)=\rho(1-\rho)$, whereas here the abandonment-adjusted value of Corollary~\ref{cor:C21:pi0} enters.

\begin{proof}[Proof of Corollary~\ref{cor:C21:pi0}]
    Two evaluations of the fundamental equation~\eqref{eq:C21:fundamental} suffice. First set $y=1$: the term $\mu x(1-y)\pi(0,0)$ vanishes, the coefficient of $P(x,1)$ factors as $-(x-1)\bigl(\lambda_1 x-(\mu+\theta_1)\bigr)$, and the right-hand side becomes $(x-1)(\mu+\theta_1)P_y(1)$. Cancelling $(x-1)$ we obtain
    \begin{equation}
        P(x,1) \;=\; \frac{(\mu+\theta_1)\,P_y(1)}{\mu+\theta_1-\lambda_1 x}
                \;=\; \frac{P_y(1)}{1-\rho_C\,x}.
        \label{eq:C21:Px1}
    \end{equation}
    This geometric marginal shows that the priority queue clears at the effective rate
    $\mu+\theta_1$. Evaluating at $x=1$ and using $P(1,1)=1-\pi_0$ gives
    \begin{equation}
        1-\pi_0 \;=\; \frac{P_y(1)}{1-\rho_C}
        \quad\Longrightarrow\quad
        P_y(1) \;=\; (1-\pi_0)(1-\rho_C).
        \label{eq:C21:P11}
    \end{equation}
    Second, set $x=1$ in Equation~\eqref{eq:C21:fundamental}. The coefficient of $P(1,y)$ collapses to $\lambda_2 y(1-y)$ and the right-hand side to $\mu(1-y)\bigl[P_y(y)-\pi(0,0)\bigr]$, yielding the following relation
    \begin{equation}
        \lambda_2\,y\,P(1,y) \;=\; \mu\bigl[P_y(y)-\pi(0,0)\bigr].
        \label{eq:C21:P1y}
    \end{equation}
    Letting $y\to1$, and using Equation~\eqref{eq:C21:P11} together with the idle-boundary balance equation~\eqref{eq:C21:idle} ---$\pi(0,0)=\lambda\pi_0/\mu$, as in Theorem~\ref{thm:gen:boundary}--- we have
    \begin{equation*}
        \lambda_2(1-\pi_0) \;=\; \mu\bigl[P_y(1)-\pi(0,0)\bigr]
        \;=\; \mu(1-\pi_0)(1-\rho_C) - \lambda\pi_0.
    \end{equation*}
    Expanding and collecting the $\pi_0$ terms,
    \begin{equation*}
        \bigl[\mu(1-\rho_C)+\lambda-\lambda_2\bigr]\pi_0 = \mu(1-\rho_C)-\lambda_2,
        \qquad\text{i.e.}\qquad
        \pi_0 = \frac{\mu(1-\rho_C)-\lambda_2}{\mu(1-\rho_C)+\lambda_1},
    \end{equation*}
    where $\lambda-\lambda_2=\lambda_1$. Substituting $\mu(1-\rho_C)=\mu(\mu+\theta_1-\lambda_1)/(\mu+\theta_1)$ and multiplying numerator and denominator by $\mu+\theta_1$, the numerator becomes
    $\mu(\mu+\theta_1-\lambda_1)-\lambda_2(\mu+\theta_1)=(\mu+\theta_1)(\mu-\lambda)+\lambda_1\theta_1$
    and the denominator $\mu(\mu+\theta_1-\lambda_1)+\lambda_1(\mu+\theta_1)=\mu(\mu+\theta_1)+\lambda_1\theta_1$,
    yielding the first identity in Equation~\eqref{eq:C21:pi0_pi00}; the second follows directly from
    $\pi(0,0)=\lambda\pi_0/\mu$. The denominator $\mu(\mu+\theta_1)+\lambda_1\theta_1$ is strictly positive, so $\pi_0>0$ is equivalent to the stability condition~\eqref{eq:C21:stability}. Finally, rewriting the computed $\pi_0$ in the renewal--reward form $\pi_0=1/\bigl(1+\lambda\,\mathbb{E}[BP]\bigr)$ of Theorem~\ref{thm:gen:boundary} identifies the whole-system mean busy period as $\mathbb{E}[BP]=(\mu+\theta_1)/\bigl[(\mu+\theta_1)(\mu-\lambda)+\lambda_1\theta_1\bigr]$ ---a by-product of the computation rather than an independent argument--- which reduces to the $M/M/1$ value $(\mu-\lambda)^{-1}$ as $\theta_1\to0^+$.
\end{proof}

\begin{remark}[The two evaluations at the singular corner $(1,1)$]\label{rem:C21:order}
At $(x,y)=(1,1)$ every term of Equation~\eqref{eq:C21:fundamental} vanishes ---$(1,1)$ is a zero of the kernel and of both right-hand-side coefficients--- so evaluating the equation at that point yields only $0=0$. Restricting the equation to the line $y=1$ and to the line $x=1$, and cancelling the vanishing linear factors $(x-1)$, respectively $(1-y)$, extracts instead two \emph{directional} first-order pieces of information at this singular corner, and they are genuinely different: Equation~\eqref{eq:C21:Px1} relates the class-$1$ marginal to $P_y(1)$, while Equation~\eqref{eq:C21:P1y} relates the class-$2$ marginal to the whole boundary function. Together with the balance equation~\eqref{eq:C21:idle} they pin down the two unknown scalars $\pi_0$ and $\pi(0,0)$. The same device applies to the other models: for Model-$A$, the two evaluations combined with $\pi(0,0)=\lambda\pi_0/\mu$ yield $\pi_0=1-\rho$ directly, without invoking Theorem~\ref{thm:A:PGF}.
\end{remark}

\subsubsection{Analytical approach for Model-\texorpdfstring{$\CH$}{C2-H}}\label{sssec:C21:analytical}

\begin{proof}[Proof of Theorem~\ref{thm:C21:PGF}]
    We again apply the Kernel Method to Equation~\eqref{eq:C21:fundamental}. Dividing the coefficient of $P(x,y)$ by $y\ne0$ and rearranging gives the quadratic equation~\eqref{eq:C21:quad}, whose roots are
    \begin{equation}
        x^{\pm}(y) = \frac{A(y) \pm \sqrt{A(y)^2 - 4\lambda_1(\mu+\theta_1)}}{2\lambda_1},
        \qquad A(y):=\lambda_1+\lambda_2(1-y)+\mu+\theta_1,
        \label{eq:C21:xstar}
    \end{equation}
    with product $x^*(y)\,x^+(y)=(\mu+\theta_1)/\lambda_1$ by Vieta's formulas. Evaluating the quadratic polynomial $f(x):=\lambda_1 x^2-A(y)x+(\mu+\theta_1)$ at $x=1$,
    \begin{equation*}
        f(1) = \lambda_1 - A(y) + (\mu+\theta_1) = -\lambda_2(1-y).
    \end{equation*}
    For $y\in[0,1)$ this is strictly negative, so $x=1$ lies between the two roots and $x^*(y)<1<x^+(y)$; at $y=1$, $x^*(1)=1$ and $x^+(1)=(\mu+\theta_1)/\lambda_1>1$ (the latter by $\rho_C<1$). Hence, $x^*(y)$ is the unique root inside the unit disc. Differentiating Equation~\eqref{eq:C21:xstar} at $y=1$, we obtain
    \begin{equation}
        \frac{d}{dy}x^*(1) = \frac{\lambda_2}{\mu+\theta_1-\lambda_1}.
        \label{eq:C21:xstar_deriv}
    \end{equation}

    Setting $x=x^*(y)$ in Equation~\eqref{eq:C21:fundamental} makes the coefficient of $P(x,y)$ vanish; since $P$ is bounded on the closed unit polydisc, the right-hand side must vanish as well,
    \begin{equation*}
        \bigl[\mu(x^*(y)-y)+\theta_1 y(x^*(y)-1)\bigr]P_y(y) = \mu x^*(y)(1-y)\,\pi(0,0),
    \end{equation*}
    and, regrouping the bracket as
    $\mu(x^*(y)-y)+\theta_1 y(x^*(y)-1)=(\mu+\theta_1 y)\,x^*(y)-y(\mu+\theta_1)$, the boundary function is
    \begin{equation}
        P_y(y) = \frac{\mu\,x^*(y)\,(1-y)\,\pi(0,0)}
                      {(\mu+\theta_1 y)\,x^*(y)-y(\mu+\theta_1)}.
        \label{eq:C21:Py}
    \end{equation}
    Write $\Delta(y):=(\mu+\theta_1 y)x^*(y)-y(\mu+\theta_1)$ for the denominator of Equation~\eqref{eq:C21:Py}; at $y=0$, $\Delta(0)=\mu\,x^*(0)$ (with $x^*(0)>0$), so $P_y(0)=\pi(0,0)$, exactly as for Model-$\BH$ at Equation~\eqref{eq:B21:Py}.
    
    Substituting Equation~\eqref{eq:C21:Py} into Equation~\eqref{eq:C21:fundamental} and writing the kernel as $-\lambda_1 y\,(x-x^*(y))(x-x^+(y))$, the right-hand side becomes
    \begin{equation*}
        \bigl[\mu(x-y)+\theta_1 y(x-1)\bigr]P_y(y) - \mu x(1-y)\,\pi(0,0)
        = \frac{\mu(1-y)\,\pi(0,0)}{\Delta(y)}\,
          \Bigl\{\bigl[\mu(x-y)+\theta_1 y(x-1)\bigr]x^*(y) - x\,\Delta(y)\Bigr\}.
    \end{equation*}
    Substituting $\Delta(y)$, the braced expression simplifies to $(\mu+\theta_1)\,y\,(x-x^*(y))$. Dividing by the factored kernel cancels $(x-x^*(y))$, yielding
    \begin{equation*}
        P(x,y) = \frac{\mu(\mu+\theta_1)(1-y)\,\pi(0,0)\,y\,(x-x^*(y))/\Delta(y)}
                      {-\lambda_1 y\,(x-x^*(y))(x-x^+(y))}
               = \frac{\mu(\mu+\theta_1)(1-y)\,\pi(0,0)}
                      {\lambda_1\,(x^+(y)-x)\,\Delta(y)},
    \end{equation*}
    which is Equation~\eqref{eq:C21:PGF} and concludes the proof.
\end{proof}
